\section{Historical Multi-concept Trial and Correction}\label{sec:historical-autonomy}
\textbf{Recorded example.} In the saved deterministic fast-learner history
(seed 3101, Bayesian-style simulator), the first day contains a diagnostic
question and a hint. Before the next recorded student turn, an
\texttt{incomplete-objective} event produces an in-app check-in: the goal still
requires assessed evidence, and the deterministic event rule selects that
permitted action. The saved opportunity, action and delivery-message identifiers
link the decision to its persisted effect and source citation. Subsequent turns
produce further learner-state revisions; the day-15 restart check retains an
identical state hash. By the end, the token-bucket concept has eight correct and
four incorrect assessments, illustrating why a delivered check-in is not itself
proof of mastery. This is one operational trace, not evidence that the check-in
caused those later outcomes. Its source is the \href{https://github.com/horiiiiii032929/digital-twin/blob/main/research/05_evaluation/governed-full-autonomy-v2-1-multi-concept-confirmation-025-results.md}{historical multi-concept result}; the sanitized trace is retained in
\path{research/06_reports/final/professor-reader-review.md}.


\subsection{Original Reported Simulator Outcomes}
For the autonomous-support question, the deterministic longitudinal study compared 72 histories in 36 matched pairs.
A pair shares the specified learner scenario and seed, with autonomous support
in one history and reactive-only support in the other. This is a comparison
within the simulator, not a randomized study of students.
It exercised 30 virtual days of tutoring and intervention logic without external
LLM calls. Simulated mastery averaged 0.365 with autonomous support and 0.332
with the reactive control; the paired difference was 0.0324, with a 95\% interval
of 0.0040--0.0617. For pair $k$, let $m_k^{A}$ and $m_k^{R}$ be final simulated mastery under
autonomous and reactive support, respectively. Each history's mastery summarizes
its simulated concepts on a zero-to-one scale. The paired mean difference is
\begin{equation}
\overline{\Delta m}=\frac{1}{36}\sum_{k=1}^{36}(m_k^{A}-m_k^{R})=0.0324.
\label{eq:mastery-difference}
\end{equation}
This is 3.24 percentage points on the simulator's scale. It is computed from
unrounded paired values, so subtracting the displayed rounded means need not
reproduce it exactly. Its interval resamples the paired differences. Both the
means and their difference follow the simulator's learning and forgetting rules;
they are not measured student learning gains.

Predictive area under the receiver operating characteristic curve (AUROC) summarizes how well the previous learner-state estimate ranks
later correct above incorrect simulated attempts. A value of 0.5 gives no ranking
advantage; 1 indicates perfect ranking. The reported 0.466, averaged over histories with both outcomes, therefore
provides little evidence of useful prediction. A simulated intervention is labelled
wasted when its recipient is already at the mastery threshold or is unreceptive.
The reported 32.9\% is an average of per-history wasted-delivery rates, rather
than a human judgment of whether a message helped. These diagnostics reveal
weaknesses that successful consent and restart checks cannot detect.


\paragraph{Later interpretation correction.}
These numbers describe the original multi-concept implementation. The subsequent
goal-scope audit found ten unsupported completed goals among 22 completions
across eight histories, even though the concept-assessment gates had passed.
The old paired mastery difference is therefore retained as a historical
simulator outcome, not evidence for the corrected lifecycle or real learning.
The \hyperref[study-H]{goal-completion regression} uses new seeds on the same exposed concepts and reports its own
paired results; it does not overwrite or rescore the historical trial.
